\documentclass[a4paper,11pt,oneside,openany]{book}
\usepackage{exam_style}

\title{ECE120: Basic Electronics Engineering Workshop --- Examination Companion \& Solved Question Papers}
\author{LPU Department of Electronics \& Communication Engineering\\Academic Editorial Board}
\date{Edition 2026--2027}

\begin{document}

\frontmatter
\begin{titlepage}
\pagecolor{LPUNavy}
\color{white}
\vspace*{1.5cm}

\begin{center}
    {\Huge \bfseries \scshape ECE-120: Basic Electronics Engineering Workshop}\\[5mm]
    {\LARGE \bfseries The Definitive University Examination Companion \& Solved Papers}\\[8mm]
    {\color{LPUOrange}\rule{0.8\textwidth}{2pt}}\\[10mm]

    \begin{tcolorbox}[enhanced, width=0.88\textwidth, colback=white, colframe=LPUOrange, boxrule=1.5pt, arc=3mm, halign=center]
        \color{DarkSlate}
        \vspace{2mm}
        {\Large\bfseries Complete Assessment Suite (Sets A \& B for CA-1, Mid-Term, CA-2, End-Term)}\\[3mm]
        {\large\bfseries Featuring Step-by-Step Mathematical Solutions \& Laboratory Rubrics}\\[3mm]
        \begin{tabular}{cc}
            \color{LPUNavy}\textbf{Full CA-1 Papers (Units 1 \& 2)} & \quad \color{LPUNavy}\textbf{Full Mid-Term Papers (Units 1, 2, 3)} \\
            \color{LPUNavy}\textbf{Full CA-2 Papers (Units 4 \& 5)} & \quad \color{LPUNavy}\textbf{Full End-Term Papers (Units 1 to 6)} \\
        \end{tabular}
        \vspace{2mm}
    \end{tcolorbox}

    \vfill

    {\large \bfseries Prepared for Lovely Professional University (LPU) Engineering Students}\\[2mm]
    {\small Course Outcomes CO1--CO6 $\bullet$ Bloom's Taxonomy L1--L4 $\bullet$ Exact Marking Rubrics}\\[10mm]

    \begin{tcolorbox}[width=0.6\textwidth, colback=LPUNavy!80!black, colframe=white, boxrule=0.5pt, halign=center]
        \color{white}
        \textbf{LPU Department of Electronics \& Communication Engineering}\\[1mm]
        {\footnotesize Academic Session 2026--2027}
    \end{tcolorbox}
    \vspace*{1cm}
\end{center}
\end{titlepage}
\nopagecolor
\newpage

\chapter*{Curriculum Blueprint \& Examination Assessment Scheme}
\addcontentsline{toc}{chapter}{Curriculum Blueprint \& Assessment Scheme}

\section*{1. Course Articulation Matrix (Course Outcomes)}
The academic syllabus of \textbf{ECE120: Basic Electronics Engineering Workshop} maps directly to six measurable Course Outcomes (CO1--CO6):
\begin{center}
\begin{tabularx}{\textwidth}{lp{9.5cm}c}
\toprule
\textbf{CO} & \textbf{Course Outcome Description} & \textbf{Bloom's Level} \\
\midrule
\textbf{CO1} & Identify passive and active electronic components, decode color codes and markings, and operate bench instrumentation (DMM, DSO, Function Generator). & L1, L2 \\
\textbf{CO2} & Apply and experimentally validate the Voltage Divider Law and Current Divider Law, accounting for meter loading and burden errors. & L3, L4 \\
\textbf{CO3} & Characterize the forward and reverse V-I conduction properties of Silicon PN junction diodes and evaluate rectifier circuits. & L2, L3 \\
\textbf{CO4} & Verify fundamental logic gate truth tables and synthesize complex logic functions using universal NAND and NOR gates. & L3, L4 \\
\textbf{CO5} & Design, wire on solderless breadboards, and evaluate combinational arithmetic circuits including Half/Full Adders and Subtractors. & L3, L4 \\
\textbf{CO6} & Program, interface, and troubleshoot an Arduino Uno microcontroller with an active Infrared (IR) proximity sensor module. & L3, L4, L5 \\
\bottomrule
\end{tabularx}
\end{center}

\section*{2. Examination Blueprint and Weightage Allocation}
\begin{center}
\begin{tabular}{lcccc}
\toprule
\textbf{Assessment Milestones} & \textbf{Syllabus Coverage} & \textbf{Exam Format} & \textbf{Duration} & \textbf{Marks} \\
\midrule
Continuous Assessment 1 (CA-1) & Units 1 \& 2 & 6 Subjective Problems (Attempt 5) & 60 min & 30 \\
Mid-Term Examination & Units 1, 2, 3 & 10 MCQs + 4 Analytical Problems & 90 min & 30 \\
Continuous Assessment 2 (CA-2) & Units 4 \& 5 & 6 Subjective Problems (Attempt 5) & 60 min & 30 \\
End-Term Examination & Units 1 to 6 & Section A (MCQs) + B (Analytical) + C (Lab Design) & 180 min & 100 \\
\bottomrule
\end{tabular}
\end{center}
\newpage

\chapter*{Quick Workshop Formula Reference Sheet}
\addcontentsline{toc}{chapter}{Quick Workshop Formula Reference Sheet}

\begin{multicols}{2}
\subsection*{Unit 1: Components \& Instruments}
\begin{itemize}
    \item Resistor: $R = (D_1 \times 10 + D_2) \times 10^M \pm \text{Tol}\,\%$
    \item Capacitor: $C = (D_1 \times 10 + D_2) \times 10^M\,\mathrm{pF}$
    \item RMS Voltage: $V_{\mathrm{rms}} = V_m / \sqrt{2} \approx 0.707 V_m$
    \item Frequency: $f = 1 / T$
    \item LED Resistor: $R_S = (V_{CC} - V_F) / I_F$
\end{itemize}

\subsection*{Unit 2: Divider Laws}
\begin{itemize}
    \item Voltage Divider: $V_x = V_{\mathrm{in}} (R_x / R_{\mathrm{total}})$
    \item 2-Resistor: $V_2 = V_{\mathrm{in}} \frac{R_2}{R_1 + R_2}$
    \item Loaded Divider: $R_{2,\mathrm{eq}} = R_2 \parallel R_L$
    \item Current Divider: $I_x = I_T (G_x / G_{\mathrm{total}})$
    \item 2-Branch: $I_1 = I_T \frac{R_2}{R_1 + R_2}$
\end{itemize}

\subsection*{Unit 3: PN Junction Diode}
\begin{itemize}
    \item Shockley: $I_D = I_s (e^{V_D / \eta V_T} - 1)$
    \item Thermal Voltage: $V_T = kT/q \approx 26\,\mathrm{mV}$ at $300\,\mathrm{K}$
    \item Dynamic Resistance: $r_d = \eta V_T / I_D = 26\,\mathrm{mV} / I_D$
    \item Static Resistance: $R_D = V_D / I_D$
    \item Half-Wave Ripple: $\gamma = 1.21$, $\eta = 40.6\%$
    \item Full-Wave Ripple: $\gamma = 0.482$, $\eta = 81.2\%$
\end{itemize}

\subsection*{Unit 4: Digital Logic \& Universality}
\begin{itemize}
    \item De Morgan: $\overline{AB} = \overline{A} + \overline{B}$
    \item De Morgan: $\overline{A+B} = \overline{A} \cdot \overline{B}$
    \item XOR: $A \oplus B = \overline{A}B + A\overline{B}$
    \item XNOR: $\overline{A \oplus B} = AB + \overline{A}\overline{B}$
    \item NAND gate count: NOT=1, AND=2, OR=3, XOR=4
\end{itemize}

\subsection*{Unit 5: Arithmetic Circuits}
\begin{itemize}
    \item Half Adder: $S = A \oplus B$, $C = AB$
    \item Full Adder Sum: $S = A \oplus B \oplus C_{\mathrm{in}}$
    \item Full Adder Carry: $C_{\mathrm{out}} = AB + C_{\mathrm{in}}(A \oplus B)$
    \item Half Subtractor: $D = A \oplus B$, $B_{\mathrm{out}} = \overline{A}B$
    \item Full Subtractor: $D = A \oplus B \oplus B_{\mathrm{in}}$
\end{itemize}

\subsection*{Unit 6: Arduino \& Sensors}
\begin{itemize}
    \item ADC Voltage: $V = \text{Reading} \times (5.0 / 1024)\,\mathrm{V}$
    \item Step size: $V_{LSB} = 4.8828\,\mathrm{mV}$
    \item Maximum Pin Current: $40\,\mathrm{mA}$ absolute limit
    \item IR Comparator (LM393): Active LOW output ($0\,\mathrm{V}$ on obstacle)
\end{itemize}
\end{multicols}
\newpage


\tableofcontents
\newpage

\mainmatter

% ==============================================================================
% SECTION 1: CONTINUOUS ASSESSMENT 1 (CA-1) --- UNITS 1 & 2
% ==============================================================================
\chapter{Continuous Assessment 1 (CA-1) Examination Suite}
\label{chap:ca1}
\noindent
Continuous Assessment 1 evaluates students on \textbf{Unit 1 (Workshop Instruments, Components, Trainer Module \& Display Devices)} and \textbf{Unit 2 (Voltage and Current Divider Laws \& Prototyping Measurements)}. Each paper contains \textbf{6 comprehensive questions} carrying \textbf{5 marks each}, and students attempt \textbf{5 questions} (Total = \textbf{25/30 Marks}, Time Allowed = \textbf{60 Minutes}).

\newpage
\section{Continuous Assessment 1 (CA-1) --- Set A}

\LPUPaperHeader{Continuous Assessment 1 (CA-1)}{A}{60 Minutes}{25 Marks}{CO1 (Components \& Instruments), CO2 (Divider Laws)}

\noindent\textbf{Instructions to Candidates:}
\begin{enumerate*}[label=\textbf{\arabic*.}]
    \item Answer any \textbf{FIVE} questions out of the six questions presented.
    \item Each question carries \textbf{5 Marks}. Clearly state governing circuit theorems.
\end{enumerate*}

\questionitem{1}{A carbon resistor has four colored bands: Brown, Black, Orange, and Gold. Determine: (a) Nominal resistance, (b) Maximum allowable resistance, (c) Minimum allowable resistance.}{5}{CO1}{L2}
\questionitem{2}{Explain the operating principle of a Digital Multimeter when measuring resistance. Why must the circuit power be turned off before connecting DMM ohmmeter test leads?}{5}{CO1}{L2}
\questionitem{3}{A voltage divider is constructed with $V_{\mathrm{in}} = 15\,\mathrm{V}$, $R_1 = 3.3\,\mathrm{k}\Omega$, and $R_2 = 6.8\,\mathrm{k}\Omega$. Calculate: (a) The unloaded output voltage across $R_2$, (b) The loaded output voltage when a load resistor $R_L = 10\,\mathrm{k}\Omega$ is connected across $R_2$.}{5}{CO2}{L3}
\questionitem{4}{State the Current Divider Law. For a parallel circuit with $I_T = 24\,\mathrm{mA}$ entering two branches with $R_1 = 4\,\mathrm{k}\Omega$ and $R_2 = 12\,\mathrm{k}\Omega$, calculate the current in each branch.}{5}{CO2}{L3}
\questionitem{5}{A student measures a sinusoidal voltage on a DSO. The peak-to-peak waveform occupies 6 vertical divisions at $2.0\,\mathrm{V/div}$, and one cycle spans 4 horizontal divisions at $0.5\,\mathrm{ms/div}$. Determine: (a) Peak voltage $V_m$, (b) RMS voltage $V_{\mathrm{rms}}$, (c) Signal frequency $f$.}{5}{CO1}{L3}
\questionitem{6}{A $5\,\mathrm{mm}$ red LED has a forward voltage drop of $V_F = 1.9\,\mathrm{V}$ at $I_F = 15\,\mathrm{mA}$. Design a current-limiting resistor to drive this LED safely from a $+5.0\,\mathrm{V}$ digital logic output. Select the nearest standard $5\%$ resistor and compute the actual LED current.}{5}{CO1}{L4}

\subsection*{Model Solutions \& Marking Rubrics (Set A)}

\begin{solutionbox}[Solution to Q.1]
(a) Brown = 1, Black = 0, Orange = $10^3$, Gold = $\pm 5\%$.\\
Nominal resistance: $R = 10 \times 10^3\,\Omega = 10\,\mathrm{k}\Omega$.\\
(b) Maximum limit: $R_{\max} = 10\,\mathrm{k}\Omega + (0.05 \times 10\,\mathrm{k}\Omega) = 10.5\,\mathrm{k}\Omega$.\\
(c) Minimum limit: $R_{\min} = 10\,\mathrm{k}\Omega - (0.05 \times 10\,\mathrm{k}\Omega) = 9.5\,\mathrm{k}\Omega$.
\end{solutionbox}
\begin{markingrubric}
\textbf{Marking Rubric (Q.1):} Decoding band values: 2 Marks $\bullet$ Calculating nominal $10\,\mathrm{k}\Omega$: 1 Mark $\bullet$ Calculating $R_{\max}$ and $R_{\min}$: 2 Marks. Total = 5 Marks.
\end{markingrubric}

\begin{solutionbox}[Solution to Q.3]
(a) Unloaded output voltage:
\[ V_{\mathrm{out,unloaded}} = V_{\mathrm{in}} \frac{R_2}{R_1 + R_2} = 15 \times \frac{6.8}{3.3 + 6.8} = 15 \times \frac{6.8}{10.1} = 10.099\,\mathrm{V} \]
(b) When $R_L = 10\,\mathrm{k}\Omega$ is attached in parallel with $R_2$:
\[ R_{2,\mathrm{eq}} = \frac{6.8 \times 10}{6.8 + 10} = \frac{68}{16.8} = 4.048\,\mathrm{k}\Omega \]
Loaded voltage:
\[ V_{\mathrm{out,loaded}} = 15 \times \frac{4.048}{3.3 + 4.048} = 15 \times \frac{4.048}{7.348} = 8.263\,\mathrm{V} \]
\end{solutionbox}
\begin{markingrubric}
\textbf{Marking Rubric (Q.3):} Correct unloaded divider formula and result ($10.10\,\mathrm{V}$): 2 Marks $\bullet$ Calculation of parallel equivalent $R_2 \parallel R_L$: 1.5 Marks $\bullet$ Loaded output computation ($8.26\,\mathrm{V}$): 1.5 Marks. Total = 5 Marks.
\end{markingrubric}
\newpage

\section{Continuous Assessment 1 (CA-1) --- Set B}

\LPUPaperHeader{Continuous Assessment 1 (CA-1)}{B}{60 Minutes}{25 Marks}{CO1 (Components \& Instruments), CO2 (Divider Laws)}

\noindent\textbf{Instructions to Candidates:}
\begin{enumerate*}[label=\textbf{\arabic*.}]
    \item Answer any \textbf{FIVE} questions out of the six questions presented.
    \item Each question carries \textbf{5 Marks}. Clearly state governing circuit theorems.
\end{enumerate*}

\questionitem{1}{A precision metal-film resistor has five color bands: Red, Violet, Black, Red, Brown. Determine: (a) Nominal resistance, (b) Tolerance percentage, (c) Upper and lower tolerance limits.}{5}{CO1}{L2}
\questionitem{2}{Differentiate between an analog multimeter and a digital multimeter (DMM) in terms of input impedance, loading effect on high-resistance circuits, and measurement accuracy.}{5}{CO1}{L2}
\questionitem{3}{A $10\,\mathrm{k}\Omega$ linear potentiometer is connected across a $+12.0\,\mathrm{V}$ regulated supply. (a) Determine the output voltage when the wiper is positioned at $35\%$ of its travel from ground. (b) If a load of $10\,\mathrm{k}\Omega$ is connected to the wiper at this position, calculate the loaded voltage.}{5}{CO2}{L3}
\questionitem{4}{A total current $I_T = 50\,\mathrm{mA}$ splits across three parallel branches with resistances $R_1 = 20\,\Omega$, $R_2 = 30\,\Omega$, and $R_3 = 60\,\Omega$. Using the conductance form of the Current Divider Law, compute the current in each branch.}{5}{CO2}{L3}
\questionitem{5}{A disc ceramic capacitor is marked with the code \textbf{473J}. (a) Decode its capacitance in picofarads, nanofarads, and microfarads. (b) State its tolerance percentage.}{5}{CO1}{L2}
\questionitem{6}{An ammeter with internal shunt resistance $R_A = 2.5\,\Omega$ is inserted to measure current in a series circuit having $V_S = 5.0\,\mathrm{V}$ and $R = 50\,\Omega$. Calculate: (a) The true circuit current without the meter, (b) The measured circuit current, (c) The percentage insertion error.}{5}{CO2}{L4}

\subsection*{Model Solutions \& Marking Rubrics (Set B)}

\begin{solutionbox}[Solution to Q.1]
(a) Red = 2, Violet = 7, Black = 0, Red = $10^2 = 100$, Brown = $\pm 1\%$.\\
Nominal resistance: $R = 270 \times 10^2\,\Omega = 27{,}000\,\Omega = 27.0\,\mathrm{k}\Omega$.\\
(b) Tolerance = $\pm 1\%$.\\
(c) Allowed range: $\pm 1\% \times 27\,\mathrm{k}\Omega = \pm 0.27\,\mathrm{k}\Omega$.\\
$R_{\max} = 27.27\,\mathrm{k}\Omega$, $R_{\min} = 26.73\,\mathrm{k}\Omega$.
\end{solutionbox}
\begin{markingrubric}
\textbf{Marking Rubric (Q.1):} 5-band decoding: 2 Marks $\bullet$ Nominal value $27\,\mathrm{k}\Omega$: 1 Mark $\bullet$ Tolerance bounds: 2 Marks. Total = 5 Marks.
\end{markingrubric}

\begin{solutionbox}[Solution to Q.4]
Conductances:
\[ G_1 = \frac{1}{20} = 0.05\,\mathrm{S}, \quad G_2 = \frac{1}{30} = 0.0333\,\mathrm{S}, \quad G_3 = \frac{1}{60} = 0.0167\,\mathrm{S} \]
Total conductance:
\[ G_{\mathrm{total}} = 0.05 + 0.0333 + 0.0167 = 0.10\,\mathrm{S} \]
Branch currents:
\[ I_1 = I_T \frac{G_1}{G_{\mathrm{total}}} = 50 \times \frac{0.05}{0.10} = 25.0\,\mathrm{mA} \]
\[ I_2 = I_T \frac{G_2}{G_{\mathrm{total}}} = 50 \times \frac{0.0333}{0.10} = 16.67\,\mathrm{mA} \]
\[ I_3 = I_T \frac{G_3}{G_{\mathrm{total}}} = 50 \times \frac{0.0167}{0.10} = 8.33\,\mathrm{mA} \]
Check: $25.0 + 16.67 + 8.33 = 50.0\,\mathrm{mA}$ (KCL verified).
\end{solutionbox}
\begin{markingrubric}
\textbf{Marking Rubric (Q.4):} Conductance calculations: 2 Marks $\bullet$ Branch current evaluations: 2.5 Marks $\bullet$ Verification of KCL sum: 0.5 Marks. Total = 5 Marks.
\end{markingrubric}
\newpage


% ==============================================================================
% SECTION 2: MID-TERM EXAMINATION SUITE --- UNITS 1, 2 & 3
% ==============================================================================
\chapter{Mid-Term University Examination Suite (Units 1, 2 \& 3)}
\label{chap:midterm}
\noindent
The Mid-Term Examination is an assessment evaluating students across \textbf{Unit 1 (Components \& Test Instruments)}, \textbf{Unit 2 (Divider Laws \& Loading Effects)}, and \textbf{Unit 3 (Semiconductor Diode Characteristics \& Rectifiers)}. Each paper consists of \textbf{10 Objective MCQs} (1 Mark each) and \textbf{4 In-Depth Analytical/Practical Design Problems} (5 Marks each) with complete step-by-step solutions (Total = \textbf{30 Marks}, Time Allowed = \textbf{90 Minutes}).

\newpage
\section{Mid-Term Examination --- Set A}

\LPUPaperHeader{Mid-Term Examination}{A}{90 Minutes}{30 Marks}{CO1, CO2, CO3 (Units 1, 2, 3)}

\subsection*{Part I: Multiple-Choice Questions (10 Marks, 1 Mark Each)}
\mcqitem{1}{Which component opposes changes in the flow of electric current?}{Capacitor}{Resistor}{Inductor}{Transformer}{1}
\mcqitem{2}{A 4-band resistor with colors Yellow, Violet, Brown, Gold has nominal value:}{$470\,\Omega$}{$4.7\,\mathrm{k}\Omega$}{$47\,\mathrm{k}\Omega$}{$47\,\Omega$}{1}
\mcqitem{3}{The input impedance of a standard modern bench Digital Multimeter in DC voltage mode is:}{$100\,\Omega$}{$10\,\mathrm{k}\Omega$}{$1\,\mathrm{M}\Omega$}{$10\,\mathrm{M}\Omega$}{1}
\mcqitem{4}{In a two-resistor voltage divider with $R_1 = R_2$, the unloaded output voltage across $R_2$ is:}{$V_{\mathrm{in}}$}{$V_{\mathrm{in}}/2$}{$2 V_{\mathrm{in}}$}{$0\,\mathrm{V}$}{1}
\mcqitem{5}{A current divider consists of $R_1 = 10\,\Omega$ and $R_2 = 40\,\Omega$. The fraction of total current flowing through $R_1$ is:}{$10/50$}{$40/50 = 0.8$}{$50/40$}{$1/2$}{1}
\mcqitem{6}{The cut-in voltage of a standard Silicon PN junction diode at room temperature is approximately:}{$0.1\,\mathrm{V}$}{$0.3\,\mathrm{V}$}{$0.7\,\mathrm{V}$}{$1.5\,\mathrm{V}$}{1}
\mcqitem{7}{The thermal voltage $V_T$ at standard room temperature ($T = 300\,\mathrm{K}$) is approximately:}{$1.0\,\mathrm{V}$}{$26\,\mathrm{mV}$}{$0.7\,\mathrm{V}$}{$50\,\mu\mathrm{V}$}{1}
\mcqitem{8}{The ripple factor of an unfiltered full-wave center-tapped rectifier is:}{$1.21$}{$0.482$}{$0.812$}{$0.05$}{1}
\mcqitem{9}{In a full-wave bridge rectifier, the Peak Inverse Voltage (PIV) rating required for each diode is:}{$V_m/2$}{$V_m$}{$2 V_m$}{$4 V_m$}{1}
\mcqitem{10}{When testing a healthy Silicon diode with a DMM in diode check mode in forward bias, the display reads:}{$0.00\,\mathrm{V}$}{$0.6\text{--}0.7\,\mathrm{V}$}{OL (Overload)}{$5.00\,\mathrm{V}$}{1}

\subsection*{Part II: Analytical \& Design Problems (20 Marks, 5 Marks Each)}
\questionitem{11}{A circuit has $V_{\mathrm{in}} = 24\,\mathrm{V}$, $R_1 = 12\,\mathrm{k}\Omega$, and $R_2 = 12\,\mathrm{k}\Omega$. Calculate the loading error percentage if an analog voltmeter with internal resistance $R_m = 24\,\mathrm{k}\Omega$ is placed across $R_2$.}{5}{CO2}{L3}
\questionitem{12}{A Silicon PN junction diode has reverse saturation current $I_s = 5\,\mathrm{nA}$ at $300\,\mathrm{K}$ ($\eta = 1$). (a) Compute forward current $I_D$ at $V_D = 0.68\,\mathrm{V}$. (b) Compute static resistance $R_D$ and dynamic resistance $r_d$ at this point.}{5}{CO3}{L3}
\questionitem{13}{Design a full-wave bridge rectifier circuit fed from a $230\,\mathrm{V}_{\mathrm{rms}}, 50\,\mathrm{Hz}$ AC mains via a $10:1$ step-down transformer. Determine: (a) Peak secondary voltage $V_m$, (b) DC output voltage $V_{dc}$, (c) PIV rating per diode.}{5}{CO3}{L4}
\questionitem{14}{A sinusoidal waveform displayed on a DSO has a period of $T = 2.5\,\mathrm{ms}$ and peak-to-peak voltage $V_{p\text{-}p} = 14\,\mathrm{V}$. (a) Calculate signal frequency. (b) Calculate RMS voltage. (c) If horizontal timebase is $0.5\,\mathrm{ms/div}$, how many divisions span one full cycle?}{5}{CO1}{L3}

\subsection*{Model Solutions \& Marking Rubrics (Mid-Term Set A)}

\begin{answerkeytable}{Part I: MCQ Answer Key}
1: C \quad 2: A \quad 3: D \quad 4: B \quad 5: B \quad 6: C \quad 7: B \quad 8: B \quad 9: B \quad 10: B
\end{answerkeytable}

\begin{solutionbox}[Solution to Q.12]
(a) Shockley equation ($V_D / \eta V_T = 0.68 / 0.026 = 26.15 \gg 1$):
\[ I_D = 5 \times 10^{-9} \times e^{26.15} = 5 \times 10^{-9} \times 2.276 \times 10^{11} = 1.138\,\mathrm{A} \]
(b) Static resistance:
\[ R_D = \frac{V_D}{I_D} = \frac{0.68\,\mathrm{V}}{1.138\,\mathrm{A}} = 0.597\,\Omega \]
Dynamic resistance:
\[ r_d = \frac{\eta V_T}{I_D} = \frac{0.026\,\mathrm{V}}{1.138\,\mathrm{A}} = 0.0228\,\Omega \]
\end{solutionbox}
\begin{markingrubric}
\textbf{Marking Rubric (Q.12):} Shockley calculation of $I_D$: 2.5 Marks $\bullet$ Static resistance $R_D$: 1 Mark $\bullet$ Dynamic resistance $r_d$: 1.5 Marks. Total = 5 Marks.
\end{markingrubric}
\newpage

\section{Mid-Term Examination --- Set B}

\LPUPaperHeader{Mid-Term Examination}{B}{90 Minutes}{30 Marks}{CO1, CO2, CO3 (Units 1, 2, 3)}

\subsection*{Part I: Multiple-Choice Questions (10 Marks, 1 Mark Each)}
\mcqitem{1}{Which passive component stores energy in an electrostatic field?}{Resistor}{Inductor}{Capacitor}{Thermistor}{1}
\mcqitem{2}{A resistor with color bands Red, Red, Orange, Gold has nominal value:}{$2.2\,\mathrm{k}\Omega$}{$22\,\mathrm{k}\Omega$}{$220\,\mathrm{k}\Omega$}{$22\,\Omega$}{1}
\mcqitem{3}{When an ammeter is connected into a circuit to measure branch current, it must be wired in:}{Parallel}{Series}{Bridge}{Star}{1}
\mcqitem{4}{The purpose of a potentiometer in an electronic workshop is:}{Store charge}{Provide an adjustable voltage divider}{Amplify small AC signals}{Step down AC mains}{1}
\mcqitem{5}{In a parallel circuit of two resistors $R_1$ and $R_2$, if $R_1 \ll R_2$, most of the incoming current flows through:}{$R_2$}{$R_1$}{Equally through both}{Neither}{1}
\mcqitem{6}{The barrier potential of a Germanium diode at room temperature is approximately:}{$0.7\,\mathrm{V}$}{$0.3\,\mathrm{V}$}{$1.2\,\mathrm{V}$}{$0.026\,\mathrm{V}$}{1}
\mcqitem{7}{As junction temperature increases, the reverse saturation current $I_s$ of a diode:}{Decreases exponentially}{Remains perfectly constant}{Doubles approximately every $10^\circ\mathrm{C}$ rise}{Drops to zero}{1}
\mcqitem{8}{The maximum theoretical rectification efficiency of a half-wave rectifier is:}{$40.6\%$}{$81.2\%$}{$50.0\%$}{$100\%$}{1}
\mcqitem{9}{In a center-tapped full-wave rectifier with peak secondary winding voltage $V_m$ across each half, the required diode PIV is:}{$V_m$}{$2 V_m$}{$4 V_m$}{$V_m / 2$}{1}
\mcqitem{10}{A shorted diode tested with a DMM in diode check mode displays:}{OL in both directions}{$\approx 0.00\,\mathrm{V}$ in both directions}{$0.7\,\mathrm{V}$ forward and OL reverse}{$10\,\mathrm{V}$}{1}

\subsection*{Part II: Analytical \& Design Problems (20 Marks, 5 Marks Each)}
\questionitem{11}{A series voltage divider consists of $V_{\mathrm{in}} = 10\,\mathrm{V}$ and two $50\,\mathrm{k}\Omega$ resistors. Calculate: (a) Unloaded output voltage across $R_2$, (b) Output voltage measured by a DMM with $10\,\mathrm{M}\Omega$ input resistance, (c) Output voltage measured by an old analog meter with $100\,\mathrm{k}\Omega$ input resistance.}{5}{CO2}{L3}
\questionitem{12}{A Silicon diode conducts $5\,\mathrm{mA}$ at $V_D = 0.65\,\mathrm{V}$ and $20\,\mathrm{mA}$ at $V_D = 0.71\,\mathrm{V}$. Calculate: (a) Dynamic resistance $r_d$ between these two operating points, (b) Static resistance at the $20\,\mathrm{mA}$ operating point.}{5}{CO3}{L3}
\questionitem{13}{A center-tapped full-wave rectifier delivers power to a $500\,\Omega$ load. If the AC voltage across each half of the secondary is $12\,\mathrm{V}_{\mathrm{rms}}$, calculate: (a) Peak voltage $V_m$, (b) DC load voltage $V_{dc}$, (c) Average DC load current $I_{dc}$, (d) Ripple voltage frequency for $50\,\mathrm{Hz}$ input.}{5}{CO3}{L4}
\questionitem{14}{A capacitor with code \textbf{104K} is placed in parallel with a resistor of $10\,\mathrm{k}\Omega$. (a) Determine the capacitance in microfarads. (b) Calculate the circuit time constant $\tau = R C$.}{5}{CO1}{L3}

\subsection*{Model Solutions \& Marking Rubrics (Mid-Term Set B)}

\begin{answerkeytable}{Part I: MCQ Answer Key}
1: C \quad 2: B \quad 3: B \quad 4: B \quad 5: B \quad 6: B \quad 7: C \quad 8: A \quad 9: B \quad 10: B
\end{answerkeytable}

\begin{solutionbox}[Solution to Q.11]
(a) Unloaded output: $V_{\mathrm{out}} = 10 \times \frac{50}{50 + 50} = 5.00\,\mathrm{V}$.\\
(b) With DMM ($R_m = 10\,\mathrm{M}\Omega = 10{,}000\,\mathrm{k}\Omega$):
\[ R_{2,\mathrm{eq}} = \frac{50 \times 10000}{50 + 10000} = 49.75\,\mathrm{k}\Omega \]
\[ V_{\mathrm{meas}} = 10 \times \frac{49.75}{50 + 49.75} = 4.987\,\mathrm{V} \quad (\text{Error } = 0.25\%) \]
(c) With analog meter ($R_m = 100\,\mathrm{k}\Omega$):
\[ R_{2,\mathrm{eq}} = \frac{50 \times 100}{50 + 100} = 33.33\,\mathrm{k}\Omega \]
\[ V_{\mathrm{meas}} = 10 \times \frac{33.33}{50 + 33.33} = 4.00\,\mathrm{V} \quad (\text{Error } = 20.0\%) \]
\end{solutionbox}
\begin{markingrubric}
\textbf{Marking Rubric (Q.11):} Unloaded voltage calculation: 1 Mark $\bullet$ DMM loading calculation: 2 Marks $\bullet$ Analog meter severe loading calculation: 2 Marks. Total = 5 Marks.
\end{markingrubric}
\newpage


% ==============================================================================
% SECTION 3: CONTINUOUS ASSESSMENT 2 (CA-2) --- UNITS 4 & 5
% ==============================================================================
\chapter{Continuous Assessment 2 (CA-2) Examination Suite}
\label{chap:ca2}
\noindent
Continuous Assessment 2 evaluates students on \textbf{Unit 4 (Digital Logic Gates \& Universal Gate Implementations)} and \textbf{Unit 5 (Combinational Arithmetic \& Logic Circuits: Half/Full Adders \& Subtractors)}. Each paper contains \textbf{6 questions} carrying \textbf{5 marks each}, with students attempting \textbf{5 questions} (Total = \textbf{25/30 Marks}, Time Allowed = \textbf{60 Minutes}).

\newpage
\section{Continuous Assessment 2 (CA-2) --- Set A}

\LPUPaperHeader{Continuous Assessment 2 (CA-2)}{A}{60 Minutes}{25 Marks}{CO4 (Logic Gates \& Universality), CO5 (Adders \& Combinational)}

\noindent\textbf{Instructions to Candidates:}
\begin{enumerate*}[label=\textbf{\arabic*.}]
    \item Answer any \textbf{FIVE} questions out of the six questions presented.
    \item Each question carries \textbf{5 Marks}. Clearly state Boolean algebra steps.
\end{enumerate*}

\questionitem{1}{Draw the pinout diagram of IC 7400 (Quad 2-Input NAND) showing VCC, GND, and the four gates. State the complete truth table.}{5}{CO4}{L2}
\questionitem{2}{Prove using De Morgan's theorems that a 2-input OR gate can be realized using three 2-input NAND gates. Draw the gate schematic.}{5}{CO4}{L3}
\questionitem{3}{Formulate the truth table of a Half Adder. Derive the minimal Boolean equations for Sum ($S$) and Carry ($C$). Draw the circuit using discrete XOR and AND gates.}{5}{CO5}{L3}
\questionitem{4}{Draw the logic diagram of a Full Adder constructed using two Half Adders and an OR gate. Prove algebraically that the Carry output is $C_{\mathrm{out}} = AB + C_{\mathrm{in}}(A \oplus B)$.}{5}{CO5}{L4}
\questionitem{5}{Synthesize a 2-input XOR gate using exactly four 2-input NAND gates. Provide the Boolean proof for the output expression.}{5}{CO4}{L4}
\questionitem{6}{Design a Half Subtractor circuit. Provide its truth table, Boolean equations for Difference ($D$) and Borrow ($B_{\mathrm{out}}$), and breadboard implementation schematic.}{5}{CO5}{L3}

\subsection*{Model Solutions \& Marking Rubrics (Set A)}

\begin{solutionbox}[Solution to Q.2]
De Morgan's Theorem states: $\overline{X \cdot Y} = \overline{X} + \overline{Y}$.\\
To implement $Y = A + B$ using NAND gates:\\
1. Invert inputs $A$ and $B$ individually using two NAND gates configured as inverters:
\[ Y_1 = \overline{A \cdot A} = \overline{A}, \quad Y_2 = \overline{B \cdot B} = \overline{B} \]
2. Feed $Y_1$ and $Y_2$ into a third NAND gate:
\[ Y = \overline{Y_1 \cdot Y_2} = \overline{\overline{A} \cdot \overline{B}} \]
Applying De Morgan's theorem to the outer inversion:
\[ Y = \overline{\overline{A}} + \overline{\overline{B}} = A + B \]
This realizes the OR operation using exactly 3 NAND gates.
\end{solutionbox}
\begin{markingrubric}
\textbf{Marking Rubric (Q.2):} Statement of De Morgan's theorem: 1 Mark $\bullet$ Algebraic derivation: 2 Marks $\bullet$ Clean schematic diagram with pin connections: 2 Marks. Total = 5 Marks.
\end{markingrubric}
\newpage

\section{Continuous Assessment 2 (CA-2) --- Set B}

\LPUPaperHeader{Continuous Assessment 2 (CA-2)}{B}{60 Minutes}{25 Marks}{CO4 (Logic Gates \& Universality), CO5 (Adders \& Combinational)}

\noindent\textbf{Instructions to Candidates:}
\begin{enumerate*}[label=\textbf{\arabic*.}]
    \item Answer any \textbf{FIVE} questions out of the six questions presented.
    \item Each question carries \textbf{5 Marks}. Clearly state Boolean algebra steps.
\end{enumerate*}

\questionitem{1}{Draw the pinout diagram of IC 7402 (Quad 2-Input NOR). Highlight why its pin mapping differs from standard 7400/7408 ICs.}{5}{CO4}{L2}
\questionitem{2}{Prove algebraically that the NOR gate is universal by demonstrating the synthesis of: (a) NOT gate, (b) OR gate, (c) AND gate using only 2-input NOR gates.}{5}{CO4}{L3}
\questionitem{3}{Formulate the complete truth table of a 1-bit Full Adder with inputs $A, B, C_{\mathrm{in}}$ and outputs $S, C_{\mathrm{out}}$. Derive the minimal sum-of-products expression for $S$.}{5}{CO5}{L3}
\questionitem{4}{Explain the architecture and carry propagation mechanism of the IC 7483 4-bit binary adder. What is the total delay if each stage has $t_{\mathrm{carry}} = 10\,\mathrm{ns}$?}{5}{CO5}{L3}
\questionitem{5}{Design a Full Subtractor circuit with inputs $A, B, B_{\mathrm{in}}$. Derive minimal Boolean expressions for Difference ($D$) and Borrow ($B_{\mathrm{out}}$).}{5}{CO5}{L4}
\questionitem{6}{What is a floating input in TTL digital circuits? Explain why floating inputs create erratic lab behavior and describe how pull-up/pull-down resistors eliminate this problem.}{5}{CO4}{L3}

\subsection*{Model Solutions \& Marking Rubrics (Set B)}

\begin{solutionbox}[Solution to Q.2]
(a) NOT Gate: Tie both inputs together: $Y = \overline{A + A} = \overline{A}$ (1 NOR gate).\\
(b) OR Gate: Invert the output of a NOR gate with a second NOR gate:
\[ Y_1 = \overline{A + B}, \quad Y = \overline{Y_1 + Y_1} = \overline{\overline{A + B}} = A + B \quad (\text{2 NOR gates}) \]
(c) AND Gate: Invert inputs individually and feed into a third NOR gate:
\[ Y_1 = \overline{A}, \quad Y_2 = \overline{B} \]
\[ Y = \overline{Y_1 + Y_2} = \overline{\overline{A} + \overline{B}} \]
By De Morgan's theorem $\overline{X + Y} = \overline{X} \cdot \overline{Y}$:\\
\[ Y = \overline{\overline{A}} \cdot \overline{\overline{B}} = A \cdot B \quad (\text{3 NOR gates}) \]
\end{solutionbox}
\begin{markingrubric}
\textbf{Marking Rubric (Q.2):} NOT realization: 1 Mark $\bullet$ OR realization: 1.5 Marks $\bullet$ AND realization with De Morgan step: 2.5 Marks. Total = 5 Marks.
\end{markingrubric}
\newpage


% ==============================================================================
% SECTION 4: END-TERM COMPREHENSIVE EXAMINATION SUITE --- UNITS 1 TO 6
% ==============================================================================
\chapter{End-Term Comprehensive University Examination Suite (Units 1 to 6)}
\label{chap:endterm}
\noindent
The End-Term Examination evaluates the entire course syllabus spanning \textbf{Units 1 through 6}. Each comprehensive examination paper is structured into three sections:
\begin{itemize}
    \item \textbf{Section A (Objective \& Short Conceptual)}: 10 Questions $\times$ 2 Marks = \textbf{20 Marks}.
    \item \textbf{Section B (Analytical Derivations \& Hardware Calculations)}: 4 Questions $\times$ 10 Marks = \textbf{40 Marks}.
    \item \textbf{Section C (Laboratory Experiment Scenario \& Microcontroller Interfacing Design)}: 2 Comprehensive Practical System Designs $\times$ 20 Marks = \textbf{40 Marks}.
\end{itemize}
Total Marks: \textbf{100 Marks} (Scaled to official university grading weights), Time Allowed: \textbf{180 Minutes (3 Hours)}.

\newpage
\section{End-Term Comprehensive Examination --- Set A}

\LPUPaperHeader{End-Term Comprehensive Examination}{A}{180 Minutes (3 Hours)}{100 Marks}{CO1, CO2, CO3, CO4, CO5, CO6 (Full Syllabus)}

\subsection*{Section A: Conceptual \& Short Analytical Questions (20 Marks, 2 Marks Each)}
\noindent Answer all \textbf{TEN} questions:
\begin{enumerate}[leftmargin=1.8em, itemsep=3pt]
    \item What is the significance of the 4th band on a carbon resistor?
    \item Explain the term 'knee voltage' of a semiconductor diode.
    \item State the dual relationship between the Voltage Divider Law and the Current Divider Law.
    \item Why must unused TTL logic inputs not be left floating on a breadboard?
    \item State the two De Morgan's theorems.
    \item Write the Boolean expression for the Sum output of a Full Adder.
    \item What is the resolution in millivolts of the 10-bit ADC in an Arduino Uno with a $5.0\,\mathrm{V}$ reference?
    \item What logic level does an active-LOW LM393 IR sensor module produce when an obstacle is detected?
    \item Calculate the nominal resistance for bands: Red, Violet, Yellow, Gold.
    \item What is the function of a flyback diode connected across an inductive relay coil?
\end{enumerate}

\subsection*{Section B: Analytical Derivations \& Hardware Design (40 Marks, 10 Marks Each)}
\noindent Answer any \textbf{FOUR} out of five questions:
\questionitem{11}{(a) Derive the Voltage Divider Law for two series resistors $R_1$ and $R_2$ across $V_{\mathrm{in}}$. (b) A $12\,\mathrm{V}$ divider with $R_1 = 15\,\mathrm{k}\Omega$ and $R_2 = 15\,\mathrm{k}\Omega$ is loaded by an external $15\,\mathrm{k}\Omega$ resistor. Compute the loaded output voltage and the percentage loading error.}{10}{CO2}{L3}
\questionitem{12}{(a) State the Shockley diode equation and define all physical terms. (b) A Silicon diode has $I_s = 2\,\mathrm{nA}$ at $300\,\mathrm{K}$ ($\eta = 1$). Calculate the diode forward current and dynamic resistance $r_d$ at $V_D = 0.65\,\mathrm{V}$. (c) Contrast static and dynamic resistance graphically.}{10}{CO3}{L3}
\questionitem{13}{(a) Prove algebraically that NAND is a universal logic gate. (b) Synthesize a 2-input XOR gate using the minimum number of 2-input NAND gates. Provide the truth table and intermediate node Boolean expressions.}{10}{CO4}{L4}
\questionitem{14}{(a) Formulate the complete truth table of a Full Adder. (b) Derive minimal Boolean expressions for Sum ($S$) and Carry Out ($C_{\mathrm{out}}$). (c) Draw the implementation schematic using two Half Adders and an OR gate, proving algebraic equivalence.}{10}{CO5}{L4}
\questionitem{15}{Compare the performance of Half-Wave, Center-Tapped Full-Wave, and Bridge Rectifiers with respect to: (a) Number of diodes, (b) DC output voltage, (c) Ripple factor, (d) PIV rating per diode, (e) Rectification efficiency.}{10}{CO3}{L3}

\subsection*{Section C: Laboratory Scenario \& Embedded Firmware Design (40 Marks, 20 Marks Each)}
\noindent Answer both comprehensive practical system design problems:
\questionitem{16}{\textbf{Industrial Automated Part Counter System Design}:\\
You are tasked with engineering a conveyor belt object counter using an Arduino Uno, an active-LOW LM393 IR sensor module, a $16\times 2$ LCD or Serial telemetry monitor, and an alarm buzzer.\\
(a) Draw the complete hardware interfacing schematic showing pin connections, power supplies, and current-limiting resistors [6 Marks].\\
(b) Explain the working principle of the LM393 comparator circuit and how the distance threshold potentiometer is calibrated [4 Marks].\\
(c) Write a robust Arduino C++ sketch incorporating software debounce to count passing items and trigger the buzzer for $500\,\mathrm{ms}$ whenever the batch count reaches 100 [10 Marks].}{20}{CO6}{L5}

\questionitem{17}{\textbf{Workshop Breadboard Prototyping \& Diagnostic Audit}:\\
During a laboratory evaluation, a student builds a combinational arithmetic circuit on a digital trainer module. However, the circuit exhibits severe anomalies: the outputs flicker erratically, the IC becomes uncomfortably hot to the touch, and the DMM voltage reading across the $+5\,\mathrm{V}$ rail sags to $2.8\,\mathrm{V}$.\\
(a) Provide a structured step-by-step diagnostic fault-tree to isolate and identify the root causes of these three symptoms [8 Marks].\\
(b) Explain the specific laboratory hazards of reversed IC power pinout ($V_{CC}$ and GND transposed) and floating CMOS inputs [6 Marks].\\
(c) Outline three mandatory workshop best-practices (power-off wiring, decoupling capacitors, common ground bonding) that prevent such hardware failures [6 Marks].}{20}{CO1, CO4}{L4}

\subsection*{Model Solutions \& Marking Rubrics (End-Term Set A)}

\begin{solutionbox}[Solution to Q.11]
(a) By KVL in series loop: $V_{\mathrm{in}} = I(R_1 + R_2) \implies I = \frac{V_{\mathrm{in}}}{R_1 + R_2}$.\\
Voltage across $R_2$: $V_{\mathrm{out}} = I R_2 = V_{\mathrm{in}} \frac{R_2}{R_1 + R_2}$.\\
(b) Numerical calculation:
Unloaded voltage: $V_{\mathrm{unloaded}} = 12 \times \frac{15}{15 + 15} = 6.00\,\mathrm{V}$.\\
When loaded by $R_L = 15\,\mathrm{k}\Omega$ across $R_2$:
\[ R_{2,\mathrm{eq}} = 15\,\mathrm{k}\Omega \parallel 15\,\mathrm{k}\Omega = 7.5\,\mathrm{k}\Omega \]
\[ V_{\mathrm{loaded}} = 12 \times \frac{7.5}{15 + 7.5} = 12 \times \frac{7.5}{22.5} = 4.00\,\mathrm{V} \]
Loading error:
\[ \%\text{ Error} = \frac{6.00 - 4.00}{6.00} \times 100\% = \frac{2.00}{6.00} \times 100\% = 33.33\% \]
\end{solutionbox}
\begin{markingrubric}
\textbf{Marking Rubric (Q.11):} Derivation of VDL: 4 Marks $\bullet$ Unloaded voltage: 2 Marks $\bullet$ Loaded voltage and error calculation: 4 Marks. Total = 10 Marks.
\end{markingrubric}

\begin{solutionbox}[Solution to Q.16 (Firmware Implementation)]
\begin{verbatim}
// Industrial Conveyor Object Counter with Software Debounce
const int IR_PIN = 2;       // IR sensor digital output (active LOW)
const int BUZZER_PIN = 8;   // Buzzer actuator output
const int LED_PIN = 13;     // Status LED indicator

volatile unsigned long totalCount = 0;
unsigned long lastTriggerTime = 0;
const unsigned long DEBOUNCE_TIME = 150; // ms lockout window

void setup() {
  pinMode(IR_PIN, INPUT_PULLUP);
  pinMode(BUZZER_PIN, OUTPUT);
  pinMode(LED_PIN, OUTPUT);
  digitalWrite(BUZZER_PIN, LOW);
  digitalWrite(LED_PIN, LOW);
  
  Serial.begin(9600);
  Serial.println("--- INDUSTRIAL CONVEYOR BATCH COUNTER ACTIVE ---");
}

void loop() {
  int sensorState = digitalRead(IR_PIN);
  
  if (sensorState == LOW) { // Obstacle breaks IR beam
    unsigned long now = millis();
    if (now - lastTriggerTime > DEBOUNCE_TIME) {
      totalCount++;
      lastTriggerTime = now;
      digitalWrite(LED_PIN, HIGH);
      
      Serial.print("Batch Counter: ");
      Serial.println(totalCount);
      
      // Batch alert at every 100 units
      if (totalCount % 100 == 0) {
        Serial.println("ALERT: Batch of 100 reached! Activating buzzer.");
        digitalWrite(BUZZER_PIN, HIGH);
        delay(500);
        digitalWrite(BUZZER_PIN, LOW);
      }
      
      delay(30);
      digitalWrite(LED_PIN, LOW);
    }
  }
}
\end{verbatim}
\end{solutionbox}
\begin{markingrubric}
\textbf{Marking Rubric (Q.16):} (a) Interfacing schematic: 6 Marks $\bullet$ (b) LM393 operational theory: 4 Marks $\bullet$ (c) Debounced firmware implementation with threshold condition: 10 Marks. Total = 20 Marks.
\end{markingrubric}
\newpage

\section{End-Term Comprehensive Examination --- Set B}

\LPUPaperHeader{End-Term Comprehensive Examination}{B}{180 Minutes (3 Hours)}{100 Marks}{CO1, CO2, CO3, CO4, CO5, CO6 (Full Syllabus)}

\subsection*{Section A: Conceptual \& Short Analytical Questions (20 Marks, 2 Marks Each)}
\noindent Answer all \textbf{TEN} questions:
\begin{enumerate}[leftmargin=1.8em, itemsep=3pt]
    \item What is the color code sequence for a $1.0\,\mathrm{k}\Omega \pm 5\%$ resistor?
    \item Define ideality factor ($\eta$) in the Shockley diode equation.
    \item Why does an ammeter have low internal resistance?
    \item State De Morgan's second theorem for three variables ($A, B, C$).
    \item How many 2-input NAND gates are required to synthesize a 2-input AND gate?
    \item Write the Carry Out equation for a Full Adder in terms of $A, B$, and $C_{\mathrm{in}}$.
    \item If the analog input to an Arduino Uno is $2.5\,\mathrm{V}$, what is the corresponding 10-bit ADC reading?
    \item State the function of the potentiometer on an IR sensor module.
    \item Differentiate between dynamic resistance and static resistance of a PN junction.
    \item What is the purpose of decoupling capacitors on digital circuit power rails?
\end{enumerate}

\subsection*{Section B: Analytical Derivations \& Hardware Design (40 Marks, 10 Marks Each)}
\noindent Answer any \textbf{FOUR} out of five questions:
\questionitem{11}{(a) Derive the Current Divider Law for a two-branch parallel network. (b) A current source of $60\,\mathrm{mA}$ feeds parallel branches of $R_1 = 30\,\Omega$ and $R_2 = 60\,\Omega$. Compute the branch currents and verify Kirchhoff's Current Law.}{10}{CO2}{L3}
\questionitem{12}{(a) Explain the forward conduction characteristic of a Silicon PN diode. (b) If $I_D = 25\,\mathrm{mA}$ at $T = 300\,\mathrm{K}$ ($\eta = 1$), determine the dynamic AC resistance $r_d$. (c) Explain why the reverse saturation current increases with temperature.}{10}{CO3}{L3}
\questionitem{13}{(a) Prove algebraically that the NOR gate is universal. (b) Realize NOT, OR, AND, and XOR operations exclusively using 2-input NOR gates.}{10}{CO4}{L4}
\questionitem{14}{(a) Design a Full Subtractor circuit. Provide its truth table. (b) Derive minimal sum-of-products Boolean expressions for Difference ($D$) and Borrow Out ($B_{\mathrm{out}}$). (c) Draw the gate-level implementation.}{10}{CO5}{L4}
\questionitem{15}{A bridge rectifier with a $1000\,\mu\mathrm{F}$ filter capacitor supplies a load of $R_L = 100\,\Omega$ from a $24\,\mathrm{V}_{\mathrm{rms}}, 50\,\mathrm{Hz}$ AC source. (a) Compute peak voltage $V_m$. (b) Calculate peak-to-peak ripple voltage $V_r = \frac{I_{dc}}{2 f C}$. (c) Determine the DC load voltage.}{10}{CO3}{L4}

\subsection*{Section C: Laboratory Scenario \& Embedded Firmware Design (40 Marks, 20 Marks Each)}
\noindent Answer both comprehensive practical system design problems:
\questionitem{16}{\textbf{Automated Touchless Dispenser System Design}:\\
You are assigned to build a touchless liquid dispenser using an Arduino Uno, an LM393 active-LOW IR proximity sensor, a $5\,\mathrm{V}$ relay module, and a miniature DC pump.\\
(a) Draw the complete wiring diagram, including the flyback protection diode across the inductive relay coil [6 Marks].\\
(b) Explain why directly driving a relay coil from an Arduino digital I/O pin exceeds maximum pin current ratings and how a transistor driver or relay module protects the ATmega328P [4 Marks].\\
(c) Write a structured Arduino C++ sketch that activates the dispenser pump for exactly $1.5\,\mathrm{seconds}$ upon detecting a hand, enforcing a mandatory $3.0\,\mathrm{second}$ lockout to prevent continuous accidental re-triggering [10 Marks].}{20}{CO6}{L5}

\questionitem{17}{\textbf{Laboratory Measurement Accuracy \& Instrument Loading Audit}:\\
An engineering workshop team measures voltages and currents in a sensitive resistive divider network ($V_{\mathrm{in}} = 10\,\mathrm{V}, R_1 = 100\,\mathrm{k}\Omega, R_2 = 100\,\mathrm{k}\Omega$) using two different instruments: an older moving-coil analog multimeter ($20\,\mathrm{k}\Omega/\mathrm{V}$ on $10\,\mathrm{V}$ range $\implies R_m = 200\,\mathrm{k}\Omega$) and a modern bench DMM ($R_m = 10\,\mathrm{M}\Omega$).\\
(a) Calculate the theoretical unloaded voltage and the voltage measured by both instruments [8 Marks].\\
(b) Compute the percentage measurement error for each meter and interpret why high circuit resistance magnifies instrument loading [6 Marks].\\
(c) Recommend two engineering solutions (instrument selection, active buffer buffering) to eliminate measurement loading errors in high-impedance circuits [6 Marks].}{20}{CO1, CO2}{L4}

\subsection*{Model Solutions \& Marking Rubrics (End-Term Set B)}

\begin{solutionbox}[Solution to Q.11]
(a) For parallel resistors $R_1$ and $R_2$, equivalent resistance is $R_p = \frac{R_1 R_2}{R_1 + R_2}$.\\
Voltage across parallel pair: $V_p = I_T R_p = I_T \left(\frac{R_1 R_2}{R_1 + R_2}\right)$.\\
Current in branch 1: $I_1 = \frac{V_p}{R_1} = I_T \left(\frac{R_2}{R_1 + R_2}\right)$.\\
Current in branch 2: $I_2 = \frac{V_p}{R_2} = I_T \left(\frac{R_1}{R_1 + R_2}\right)$.\\
(b) Numerical calculation:
\[ I_1 = 60\,\mathrm{mA} \times \frac{60}{30 + 60} = 60 \times \frac{60}{90} = 40.0\,\mathrm{mA} \]
\[ I_2 = 60\,\mathrm{mA} \times \frac{30}{30 + 60} = 60 \times \frac{30}{90} = 20.0\,\mathrm{mA} \]
Sum: $I_1 + I_2 = 40.0 + 20.0 = 60.0\,\mathrm{mA} = I_T$ (KCL verified).
\end{solutionbox}
\begin{markingrubric}
\textbf{Marking Rubric (Q.11):} CDL algebraic derivation: 4 Marks $\bullet$ Branch current calculations: 4 Marks $\bullet$ KCL check: 2 Marks. Total = 10 Marks.
\end{markingrubric}

\begin{solutionbox}[Solution to Q.16 (Firmware Implementation)]
\begin{verbatim}
// Touchless Dispenser Controller with Hold-off Lockout
const int IR_PIN = 2;          // Active LOW IR sensor input
const int RELAY_PIN = 7;       // Relay control output pin
const int STATUS_LED = 13;     // Status LED indicator

const unsigned long DISPENSE_TIME = 1500; // 1.5 seconds pump active
const unsigned long LOCKOUT_TIME = 3000;  // 3.0 seconds cooldown period

unsigned long dispenseStartTime = 0;
unsigned long dispenseEndTime = 0;
bool isDispensing = false;

void setup() {
  pinMode(IR_PIN, INPUT_PULLUP);
  pinMode(RELAY_PIN, OUTPUT);
  pinMode(STATUS_LED, OUTPUT);
  digitalWrite(RELAY_PIN, LOW); // Relay OFF
  digitalWrite(STATUS_LED, LOW);
  
  Serial.begin(9600);
  Serial.println("--- TOUCHLESS DISPENSER INITIALIZED ---");
}

void loop() {
  unsigned long currentMillis = millis();
  int sensorVal = digitalRead(IR_PIN);
  
  if (!isDispensing) {
    // Check if cooldown lockout has elapsed
    if (currentMillis - dispenseEndTime > LOCKOUT_TIME) {
      if (sensorVal == LOW) { // Hand detected
        Serial.println("Hand Detected: Activating Dispenser Pump...");
        isDispensing = true;
        dispenseStartTime = currentMillis;
        digitalWrite(RELAY_PIN, HIGH); // Pump ON
        digitalWrite(STATUS_LED, HIGH);
      }
    }
  } else {
    // Check if dispense duration has completed
    if (currentMillis - dispenseStartTime >= DISPENSE_TIME) {
      Serial.println("Dispense Complete. Entering Cooldown Period...");
      digitalWrite(RELAY_PIN, LOW); // Pump OFF
      digitalWrite(STATUS_LED, LOW);
      isDispensing = false;
      dispenseEndTime = currentMillis;
    }
  }
}
\end{verbatim}
\end{solutionbox}
\begin{markingrubric}
\textbf{Marking Rubric (Q.16):} (a) Schematic with flyback diode: 6 Marks $\bullet$ (b) ATmega328P current limit explanation ($40\,\mathrm{mA}$ max): 4 Marks $\bullet$ (c) Non-blocking timing firmware with state machine logic: 10 Marks. Total = 20 Marks.
\end{markingrubric}
\newpage


\end{document}
